\section{Introduction}
\label{sec:intro}

% P1 -- the problem and why it matters (no implementation detail).
\note{P1: Predicting transcriptional response to unseen perturbations is central to drug discovery and
virtual-cell efforts; networks (PPI, signalling, KGs) are the default way to inject biology.}
% Cite: L1000 (Subramanian 2017), Tahoe-100M (Zhang 2025), DeepCE (Pham 2021), GSNN (Evans 2024),
% GEARS, PDGrapher, TxPert, PrimeKG-based models (Passigan 2025).

% P2 -- the gap: published network ablations use controls that cannot separate degree from wiring.
\note{P2: The evidence that networks help rests on comparisons with uniform random graphs or no graph.
Those controls change the degree distribution, and degree alone is a strong predictor in gene networks
(Gillis \& Pavlidis 2011); pathway masks match sparsity-matched random masks (Caranzano 2025); random graphs
work for expression feature selection (Bertin 2019). GSNN's +24\% over a random graph is a concrete case:
its control redraws all three edge sets without preserving degree.}

% P3 -- our approach.
\note{P3: We separate the two with (i) a control ladder applied to any model's graph and (ii) a model-free
information test asking what a network can carry about responses at all, so a null result is not
blamed on the training.}

% P4 -- headline results (strongest first).
\note{P4: Degree, not wiring: target degree predicts response magnitude in every curated network and
dataset; in models, real $\approx$ rewired everywhere we test; directed regulatory networks are the
exception at the gene level. Positive side: target identity, morphology and one measured response carry
the transferable information for unseen drugs.}

% Contributions list (keep to 3-4 items, each one sentence, each tied to a figure).
\paragraph{Contributions.}
\begin{itemize}
  \item A control ladder and a model-free information test for network priors, released as a
        model-agnostic toolkit (\cref{sec:framework}; Fig.~1).
  \item A survey of eight networks across L1000 and Tahoe-100M showing that the response information in
        undirected interaction networks is mostly target degree, with a configuration-model argument for
        why propagation over a rewired graph retains it (\cref{sec:info}; Fig.~2).
  \item Audits of published models (DeepCE, GSNN, GEARS, PDGrapher) and of our own model families under
        the ladder (\cref{sec:audit}; Fig.~3).
  \item An information-tier analysis identifying what does transfer to unseen compounds, and directed
        regulatory networks as the network signal worth modelling (\cref{sec:tiers}; Fig.~4).
\end{itemize}
